%! Solving Polynomial Equations by Factoring — Unit 1, Lesson 1.6 — Slides (Beamer)
\documentclass[aspectratio=169, 11pt]{beamer}
\usepackage{atda-beamer}   % provides \royalheader{} and \sectionlabel[color]{}
\usepackage{pgfplots}      % atda-beamer loads tikz but NOT pgfplots; needed for the graph frame
\pgfplotsset{compat=1.18}

\begin{document}

% TITLE SLIDE (CourseName is not defined in beamer, so the name is literal)
{
\setbeamercolor{background canvas}{bg=royal}
\begin{frame}[plain]
  \vspace{2.8cm}\hspace{0.55cm}%
  \begin{minipage}{0.9\textwidth}
    {\color{goldacc}\bfseries\small LESSON 1.6}\\[0.5cm]
    {\color{white}\bfseries\LARGE Solving Polynomial Equations by Factoring}\\[1.2cm]
    {\color{mistmid}\small Algebra, Trigonometry, and Data Analysis~$\cdot$~Unit 1}
  \end{minipage}
\end{frame}
}

% ── HOOK ──────────────────────────────────────────────────────────────────────
\begin{frame}[t, plain]
  \royalheader{The T-Shirt Launcher}
  \sectionlabel[goldacc]{hook --- the crew needs to know \emph{when}}
  Fired from the top of the bleachers, $48$ feet above the gym floor:
  \[ h(t) = -16t^2+32t+48 \quad \text{feet.} \]
  \begin{block}{When does the shirt hit the floor?}
    That is the question $h(t)=0$. You already have every tool --- last night's item 13(b) ran the
    whole method once, before it had a name.
  \end{block}
\end{frame}

% ── THE WARM-UP IS THE LESSON ─────────────────────────────────────────────────
\begin{frame}[t, plain]
  \royalheader{Zero Is the Only Number That Tells You Anything}
  \sectionlabel{warm-up item 3, said out loud}
  \[ ab = 12 \;\Longrightarrow\; a \text{ could be } 1, \; 2, \; 3, \; 12, \; \tfrac12, \dots \]
  \[ ab = 0 \;\Longrightarrow\; {\color{goldacc}a = 0 \;\text{ or }\; b = 0} \]
  \begin{block}{The zero product property}
    If a product of factors is $0$, at least one factor \textbf{is} $0$. It works for any number of
    factors --- and for \textbf{no other number} on the right-hand side.
  \end{block}
\end{frame}

% ── ALREADY FACTORED ──────────────────────────────────────────────────────────
\begin{frame}[t, plain]
  \royalheader{A Factored Equation Is Almost Solved}
  \sectionlabel{set each factor equal to zero}
  \vspace{0.15cm}
  \renewcommand{\arraystretch}{1.5}
  \begin{tabularx}{\textwidth}{@{}p{3.6cm} p{4.6cm} X@{}}
    \textbf{Equation} & \textbf{Each factor $=0$} & \textbf{Solutions} \\
    \hline
    $(x-4)(x+7)=0$  & $x-4=0$; \ $x+7=0$  & $x=4$, \ $x=-7$ \\
    $3x(x-5)=0$     & $x=0$; \ $x-5=0$    & $x=0$, \ $x=5$ \\
    $(2x-1)(x+6)=0$ & $2x-1=0$; \ $x+6=0$ & $x=\tfrac12$, \ $x=-6$ \\
    $(x-3)^2=0$     & $x-3=0$ (twice)     & $x=3$ \ (double root) \\
  \end{tabularx}
  \vspace{0.3cm}

  The constant $3$ is never zero --- but the factor $x$ is, and $x=0$ is a solution you must keep.
\end{frame}

% ── THE METHOD ────────────────────────────────────────────────────────────────
\begin{frame}[t, plain]
  \royalheader{The Method --- Standard Form First}
  \sectionlabel{the property only works against zero}
  \begin{enumerate}
    \setlength{\itemsep}{0.2cm}
    \item \textbf{Set the equation equal to zero.}
    \item \textbf{Factor completely} --- the order from Lessons 1.3--1.5, GCF first.
    \item \textbf{Set each factor equal to zero} and solve.
    \item \textbf{Check} each solution in the original equation.
  \end{enumerate}
  \vspace{0.3cm}
  {\color{goldacc}\bfseries $x^2+5x=24$ is not ready.} \quad
  $x^2+5x-24=0 \;\Rightarrow\; (x+8)(x-3)=0 \;\Rightarrow\; x=-8, \; x=3$.
\end{frame}

% ── THE TRAP ──────────────────────────────────────────────────────────────────
\begin{frame}[t, plain]
  \royalheader{Never Divide Both Sides by a Variable}
  \sectionlabel[goldacc]{the trap of the day}
  \[ 2x^2 = 8x \]
  Divide by $x$ and you get $2x=8$, so $x=4$. True --- and \textbf{incomplete}.
  \begin{block}{Move it instead}
    \[ 2x^2-8x=0 \;\Rightarrow\; 2x(x-4)=0 \;\Rightarrow\;
       {\color{goldacc}x=0} \;\text{ or }\; x=4 \]
    Dividing by $x$ quietly assumes $x \ne 0$ --- and here $x=0$ is a solution.
  \end{block}
\end{frame}

% ── DEGREE 3 ──────────────────────────────────────────────────────────────────
\begin{frame}[t, plain]
  \royalheader{Degree Three --- Count Before You Solve}
  \sectionlabel{\texttt{A2.EI.6a--b} --- how many, and what kind}
  \[
  \begin{aligned}
    x^3-9x &= 0 \\
    x(x+3)(x-3) &= 0 \\
    x &= 0, \quad x = -3, \quad x = 3
  \end{aligned}
  \]
  \begin{block}{The degree is the count}
    Degree $3$ $\Rightarrow$ \textbf{three} solutions, counting a double root twice. And the GCF
    step is what produced $x=0$ --- pulling out $x$ is never optional.
  \end{block}
\end{frame}

% ── REAL VS IMAGINARY ─────────────────────────────────────────────────────────
\begin{frame}[t, plain]
  \royalheader{When a Solution Is Not a Real Number}
  \sectionlabel[goldacc]{Lesson 1.5's prime sum of squares, cashed in}
  \[
  \begin{aligned}
    x^3+4x &= 0 \\
    x\left(x^2+4\right) &= 0 \\
    x = 0 \quad&\text{or}\quad x^2 = -4
  \end{aligned}
  \]
  \begin{block}{No real number squares to a negative}
    So $x=0$ is the only \textbf{real} solution. Degree $3$ still means three solutions --- the
    other two are \textbf{imaginary}. Count them and name the type.
  \end{block}
\end{frame}

% ── BACKWARDS ─────────────────────────────────────────────────────────────────
\begin{frame}[t, plain]
  \royalheader{Running the Method Backwards}
  \sectionlabel{\texttt{A2.EI.6a} --- from solutions to an equation}
  A solution $x=r$ comes from the factor $(x-r)$. One factor per solution:
  \vspace{0.15cm}
  \renewcommand{\arraystretch}{1.5}
  \begin{tabularx}{\textwidth}{@{}p{4.4cm} X@{}}
    \textbf{Solutions wanted} & \textbf{Factored equation} \\
    \hline
    $5$ and $-3$          & $(x-5)(x+3)=0$ \\
    $0$ and $7$           & $x(x-7)=0$ \\
    $-2$, $1$, and $4$    & $(x+2)(x-1)(x-4)=0$ \\
    $0$, $-3$, and $3$    & $x(x+3)(x-3)=0$ \\
  \end{tabularx}
\end{frame}

% ── GRAPH ─────────────────────────────────────────────────────────────────────
\begin{frame}[t, plain]
  \royalheader{A Graph Hands You the Same Information}
  \sectionlabel{an $x$-intercept \emph{is} a solution}
  \begin{columns}[T]
    \begin{column}{0.52\textwidth}
      \vspace{-0.2cm}
      \begin{tikzpicture}
      \begin{axis}[width=6.4cm, height=4.6cm,
          axis lines=middle, xlabel=$x$, ylabel=$y$,
          xmin=-3.4, xmax=4.4, ymin=-11, ymax=11,
          xtick={-3,-2,-1,1,2,3,4}, ytick={-8,-4,4,8},
          tick label style={font=\tiny}, label style={font=\tiny},
          samples=101, domain=-2.6:3.5]
        \addplot[goldacc, very thick, smooth] {x^3 - x^2 - 6*x};
      \end{axis}
      \end{tikzpicture}
    \end{column}
    \begin{column}{0.44\textwidth}
      Where the curve crosses the $x$-axis, the output is $0$ --- so each $x$-intercept is a
      \textbf{zero of the function}.
      \vspace{0.25cm}

      Intercepts $-2$, $0$, $3$:
      \[ x(x+2)(x-3)=0 \]
      Any nonzero multiple works too.
    \end{column}
  \end{columns}
\end{frame}

% ── ACTIVITY LAUNCH ───────────────────────────────────────────────────────────
\begin{frame}[t, plain]
  \royalheader{Group Activity: The Model Rocket}
  \sectionlabel{groups of three --- everyone starts at Tier R}
  Launched from the field-house roof, $128$ feet up:
  \[ h(t) = -16t^2+32t+128 \quad \text{feet.} \]
  \begin{block}{File the recovery window}
    When is the rocket back on the ground? Then: the club claims it passes $128$ feet
    \textbf{twice}. Are they right?
  \end{block}
\end{frame}

% ── CLOSE ─────────────────────────────────────────────────────────────────────
\begin{frame}[t, plain]
  \royalheader{Exit Ticket}
  \sectionlabel[goldacc]{notes closed --- 5 minutes}
  \begin{enumerate}
    \setlength{\itemsep}{0.35cm}
    \item Solve: $x^2-3x-10=0$.
    \item Solve: $x^3-16x=0$. How many solutions, and what kind?
    \item A student solves $3x^2=12x$ by dividing by $3x$ and gets $x=4$. Correct? Complete?
  \end{enumerate}
  \vspace{0.3cm}
  {\color{royal}\bfseries Next: Lesson 1.7 --- Simplifying Rational Expressions.}
\end{frame}

\end{document}
